\chapter{Peak fitting}\label{Chap:peakfitting}

There are two ways that GSAS-II can fit a diffraction pattern. Most of this book is concerned with Rietveld analysis, where peak positions are generated from the structure factors generated from a space group, lattice parameters and atom positions/ADPs for one or more crystallographic phases. In the Pawley and Le Bail fitting options (see section \ref{ref:LeBail}), the reflection intensities are optimized directly, rather than being determined from the atom positions/ADPs. The second mechanism uses discrete peaks that are placed at the locations that you choose, independent from any lattice constraints. The position, intensity and, optionally, peak widths for each peak are varied independently. We call this "peak fitting" to distinguish it from the full-profile Rietveld fitting process. 

There are a number of reasons why peak fitting may be used. The capability was initially developed as part of the indexing process, where one assigns a lattice to a set of peaks. Use of peak fitting provides much more accurate positions than other methods and can uncover where two peaks are closely spaced -- the inclusion of a single peak at a position intermediate between the two actual peaks can derail the fitting process. However, there are other reasons why peak fitting can be useful. It can be used to determine the intensity of peaks, where that is the only quantity needed to be determined perhaps for phase abundance or some \textit{ad hoc} property is being characterized, possibly where lattice parameters are not known, or where only a narrow data range has been measured. It can be useful for identification of an unknown phase (see \ref{ref:extrapeak}) or for learning about peak widths in the absence of the restrictions of a model. 

It should be noted that peak fitting uses the same background fitting tools described in Chapter \ref{Chap:background}. Background parameters will be optimized in the peak fitting process when their refinement flags are on. The instrumental profile terms described in Chapter \ref{Chap:profile} may also be used, depending on the peak width mode selected (discussed further in section \ref{ref:widthmode}) and the instrumental terms (located in the histogram's Instrumental Parameters data tree child entry) can optionally be fit as well. 

To perform peak fitting, select the Peak List data tree child item for the histogram of interest. This will present a table of peaks, that will initially be empty. Each peaks is described with four parameters and each of those four parameters has a refinement flag to determine if that parameter will be optimized when peaks are refined. 

\section{Peak Parameters Definitions}\label{ref:sigmadef}

The four peak parameters are "position" with the peak position as either $2\theta$ (degrees)
or TOF ($\mu$sec), "intensity" which is the integrated intensity for the peak on the intensity scale of the diffraction pattern. The "sigma${}^2$" and "gamma" values are the Gaussian and Lorentzian peak widths, respectively. 
Note that the Gaussian full-width at half-maximum ($FWHM_G$)  is related to "sigma${}^2$" ($\sigma^2$), which is the peak variance, by: 

$$FWHM_G = \sqrt{8 * \ln 2 * \sigma^ 2}$$

while the Lorentzian full-width at half-maximum ($FWHM_L$) is just "gamma" ($\Gamma$):

$$FWHM_L = \Gamma$$

The $sigma^2$  values are in units of centidegrees$^2$ (centidegrees are degrees $\times 100$) for CW data or microseconds${}^2$ for TOF data; 
$\Gamma$ has units of centidegrees or microseconds.  

\section{Creating a Peak List}
Peaks are added to the peak list by clicking on a data point in the pattern. Note that when the plot is in zoom or pan mode (selected by the button with perpendicular arrows or the magnifying glass at the lower left), mouse clicks are not available to GSAS-II and peaks are not added. When a peak is added, a vertical line is drawn at the peak position. Right-clicking on a peak deletes it. Peaks can be moved by dragging them -- press the mouse button down at the location of  the vertical line and move the mouse to the new location for the peak while holding the mouse button down. Note that to drag a peak, click on a part of the line away from any data points. If you click on a point, you will add an additional peak. 

If a peak is selected in the table in the data window, which can be done by double-clicking on the number to the left of the peak values, the corresponding vertical line is highlighted. The "d" and "u" keyboard keys, when pressed with the plot window active, will  move successively through the peak table, selecting each peak in turn (downwards with "d" and upwards with "u").  

Peaks may added by reading them from a file. Use the "Peak List"/"Load Peaks..." menu command to read peaks from a file saved with the "Peak List"/"Save Peaks..." menu command. Do not use the  Import/"Read powder pattern peaks" menu command. That is used for a different type of histogram that only has peaks, and does not have an associated powder pattern. 

\section{Peak Width modes}\label{ref:widthmode}
There are three ways that peak widths can be managed. 
\begin{enumerate}
    \item If a peak width ($\sigma^2$ or $\Gamma$) has its refine flag turned on, then that value in the table will be used for an initial peak fit and will be refined if a peak fit is performed. 
    \item In the default mode of operation, any unrefined $\sigma^2$ or $\Gamma$ values will be set from the histogram instrumental parameters, (U, V \& W for $\sigma^2$ and X, Y \& Z for $\Gamma$). When this is done, it is possible to refine those instrumental parameters, but make sure that there are enough peaks with widths set in that manner to over-define all the instrumental terms that will be refined.
    \item Should you want to fix the $\sigma^2$ and $\Gamma$ values to specific values and not refine them, there is a mode that can be used to prevent changing peak widths to the values from the There is menu command setting,  "Peak Fitting"/"Gen unvaried widths" that is enabled (checked) by default. If this is unchecked, then the values of $\sigma^2$ and $\Gamma$ will not be changed, even when not refined. 
\end{enumerate}

My recommendation is that the Gaussian broadening,  $\sigma^2$, be generated from the instrumental values for U, V \& W. If the instrument performance has been calibrated, these values do not need to be varied. If the instrument has not had its instrumental broadening calibrated, the U, V \& W values may be refined here and these provide reasonable starting values for a calibration.  These Gaussian values will almost never be affected by sample phenomena, so they should not be affected by any type of anisotropic broadening. 

The $\Gamma$ values will almost always arise from sample broadening and may vary from peak to peak. This can be because adjacent peaks may arise from different phases or there may be anisotropic crystallite or size broadening. None of this can be modeled with refinement of the X \& Y terms. The Z term is not associated with any physical process and is not expected to be needed, but is provided to allow more freedom in Lorentzian widths generation, should that ever be needed. Thus, it does make sense to vary the individual gamma terms for each peak. 

\section{The Peak Fitting Process}\label{ref:peakparameterorder}
The process of fitting peaks is usually done in a few steps. There are times where I will refine peak intensities, positions and widths all together, usually I will follow the steps outlined here. 

Note that a peak fit refinement is performed using the "Peak Fitting"/Peakfit menu command. (The "Peak Fitting"/"Peakfit one cycle" menu command does the same thing, but stops after a single least-squares cycle.) Note that if any background terms (Background tree item) are set to be fit, they will also be refined when the Peakfit is run. Likewise, if the refine flags are set for U, V, W, X, Y and/or Z on the Instrument Parameters tree item), these parameters will also be refined. Initially you will likely want to check that the Instrument Parameters refine flags are all turned off. You may want to use the tools discussed in section \ref{ref:backgroundfit}. At least in the early stages of peak fitting, you either should fix the background or else have simple background settings where only a few background terms are fit. This is particularly true if the data range is relatively small. Do consider the upper and lower data limits, as even if you only supply peaks in a small section of the pattern, if the limits cover the entire data range, then the background will be fit over that entire range. The background over a small data region is likely to only need a few terms for fitting. 

Initially peaks are located and then the peak intensities are refined (that flag is turned on by default as new peaks are created.) The fit is done using the "Peak Fitting"/Peakfit menu command. Once the peak intensities have been fit, look for any peaks with negative or approximately zero intensity. You may want to delete these peaks, or if you feel that there is actually peak in that location and will want to put it back into the refinement at a later time, you can reset the intensity manually and and remove the intensity refinement flag for that peak In the next step, you will refine peak positions, but the position of a peak with negligible intensity has no impact on the fit and the refinement where the position value is fit unlikely to progress. Likewise, when there are two closely spaced peaks, which may appear to be a single extra wide peak or a weaker shoulder on the edge of a strong peak, it may be hard to refine the positions of both peaks. You may be able to refine both peak positions later in the process, when the peak widths are well defined, but initially one should refine peak positions only where there is a clear peak maximum visible. Select the position refinement flag for these clearly defined peak positions. It can be helpful to double-click on a row label in the table, as that will highlight the peak location in the plot. Once the peaks to have positions varied have been identified, again use the "Peak Fitting"/Peakfit menu command. 

At this point, one should consider the quality of the background fit and how well the peak widths are matched. For background fitting, all the tools discussed in section \ref{ref:backgroundfit} are available, but it is also possible to simply turn on and off the fitting of terms or change the number of terms that are fit as part of the peak fitting process. When it comes to fitting the peak profile, there is a decision to be made: do you trust the U, V \& W terms supplied in the instrument parameters file? If you do believe the calibration was done well, you do not want to refine the U, V \& W terms or for that matter the $\sigma^2$ term for any peaks. If you do not have accurate calibration for your data, you may want to refine the Gaussian widths. If you are fitting 5 or more peaks, you can consider fitting the  U, V \& W terms, but fitting a three-parameter function with 4 or fewer peaks is poor idea. In that case, refine the $\sigma^2$ widths directly. In addition, unless the pattern is very well fit at this stage, include Lorentzian broadening by fitting the $\Gamma$ terms for each peak. You may not be able to refine the individual $\Gamma$ terms for smaller peaks that are shoulders on larger peaks. 

\section{Extra Peak mode}\label{ref:extrapeak}
If you have performed a refinement, but find that there are additional peaks in the pattern that are not being indexed, you will likely want to identify the source for these peaks. 

One possibility is that the symmetry of the material has been lowered. GSAS-II provides a fairly simple way to look for that using the visualization in the Unit Cells List data tree entry. This is discussed in section \ref{ref:supercellindexing}.

More commonly, there is an unexpected phase present. In the case of neutron scattering, if it is also possible that the peaks are from magnetic scattering. To help identify an unexpected phase or to find the k-vector for the expanded magnetic cell, you will want to have accurate positions for the peaks that are not being indexed in the pattern. For this one can fit peaks to the difference plot, where the these extra peaks should be amongst the most promising features. This is done using "Extra Peak mode" which is enabled, either by clicking on the button labeled "Switch to Extra Peak mode" at the top of the data window, or set turn on the checkmark  setting for the "Add impurity/subgrp/magnetic peaks." The two options do exactly the same thing. 


\section{Sequential Peak Fitting}\label{ref:seqpeak}

(more work needed: \ref{needs more work})